\documentclass[12pt, a4paper]{article}

% ── Packages ──────────────────────────────────────────────────────
\usepackage[utf8]{inputenc}
\usepackage[T1]{fontenc}
\usepackage{amsmath, amssymb, amsthm}
\usepackage{mathtools}
\usepackage{enumitem}
\usepackage{tikz}
\usetikzlibrary{calc, positioning, arrows.meta, fit, patterns, decorations.pathreplacing}
\usepackage{geometry}
\geometry{margin=2.4cm}
\usepackage{xcolor}
\usepackage{tcolorbox}
\tcbuselibrary{theorems, skins, breakable}
\usepackage{booktabs}
\usepackage{hyperref}
\hypersetup{colorlinks=true, linkcolor=blue!60!black, urlcolor=blue!60!black}

\setlength{\parindent}{0pt}
\setlength{\parskip}{6pt}

% ── Colours ───────────────────────────────────────────────────────
\definecolor{setA}{RGB}{66, 133, 244}   % blue
\definecolor{setB}{RGB}{234, 67, 53}    % red
\definecolor{setC}{RGB}{251, 188, 4}    % yellow
\definecolor{theoryBg}{RGB}{232, 245, 253}
\definecolor{theoryFrame}{RGB}{66, 133, 244}
\definecolor{stmtBg}{RGB}{255, 243, 224}
\definecolor{stmtFrame}{RGB}{230, 126, 34}
\definecolor{ansBg}{RGB}{232, 250, 232}
\definecolor{ansFrame}{RGB}{39, 174, 96}

% ── Environments ──────────────────────────────────────────────────
\newtcolorbox{theorybox}[1][]{%
  colback=theoryBg, colframe=theoryFrame,
  fonttitle=\bfseries, title={\small Theory Recap: #1},
  breakable, enhanced, arc=2mm,
  left=5pt, right=5pt, top=5pt, bottom=5pt
}

\newtcolorbox{problembox}{%
  colback=stmtBg, colframe=stmtFrame,
  fonttitle=\bfseries, title={\small Problem Statement},
  breakable, enhanced, arc=2mm,
  left=5pt, right=5pt, top=5pt, bottom=5pt
}

\newtcolorbox{answerbox}{%
  colback=ansBg, colframe=ansFrame,
  arc=2mm, left=5pt, right=5pt, top=3pt, bottom=3pt
}

\newcommand{\Z}{\mathbb{Z}}
\newcommand{\R}{\mathbb{R}}
\newcommand{\N}{\mathbb{N}}
\newcommand{\powset}[1]{2^{#1}}
\newcommand{\comp}[1]{\overline{#1}}
\newcommand{\symdiff}{\mathbin{\triangle}}

% ── Title ─────────────────────────────────────────────────────────
\title{\textbf{Practice 5 --- Combinatorics (Basic Rules)}\\[4pt]
       \large Solutions to Selected Problems\\[2pt]
       \normalsize \S1: Problems 3, 5, 10 \quad \S2: Problems 5, 10, 12}
\author{Discrete Mathematics --- QYSU}
\date{}

\begin{document}
\maketitle
\tableofcontents
\newpage

% ══════════════════════════════════════════════════════════════════
\part*{Section 1: Fundamental Rules (Sum and Product Rules)}
\addcontentsline{toc}{part}{Section 1: Fundamental Rules}
% ══════════════════════════════════════════════════════════════════

% ══════════════════════════════════════════════════════════════════
\section{Problem 3 --- Choosing Two Books of Different Subjects}
% ══════════════════════════════════════════════════════════════════

\begin{theorybox}[Sum and Product Rules]
The two most fundamental counting principles:

\begin{itemize}[nosep]
  \item \textbf{Sum Rule (Rule of Alternatives):}
    If a task can be performed in one of $k$ mutually exclusive ways,
    and way $i$ can be done in $n_i$ ways, then the total number of
    ways is $n_1 + n_2 + \cdots + n_k$.

  \item \textbf{Product Rule (Rule of Sequential Choices):}
    If a task consists of $k$ sequential steps, where step $i$ can be
    performed in $n_i$ ways (independently of the choices in other steps),
    then the total number of ways is $n_1 \cdot n_2 \cdots n_k$.
\end{itemize}

\medskip
\textbf{Combined usage.}\;
When a task can be decomposed into several \emph{cases} (sum rule),
each of which involves \emph{sequential choices} (product rule),
we compute the count for each case using the product rule, then
add the results.
\end{theorybox}

\begin{problembox}
In how many ways can you choose two books of different subjects
if there are 15 books on Informatics, 12 on Mathematics, and 10 on Physics?
\end{problembox}

\subsection*{Solution}

Since the two chosen books must be on \emph{different} subjects,
we enumerate all possible pairs of subjects.
There are three cases (sum rule), and within each case we apply
the product rule:

\begin{center}
\renewcommand{\arraystretch}{1.4}
\begin{tabular}{lccc}
\toprule
\textbf{Case} & \textbf{Subject 1} & \textbf{Subject 2} & \textbf{Count} \\
\midrule
Informatics + Mathematics & 15 choices & 12 choices & $15 \cdot 12 = 180$ \\
Informatics + Physics     & 15 choices & 10 choices & $15 \cdot 10 = 150$ \\
Mathematics + Physics     & 12 choices & 10 choices & $12 \cdot 10 = 120$ \\
\bottomrule
\end{tabular}
\end{center}

By the sum rule (the three cases are mutually exclusive), the total is:
\[
  15 \cdot 12 \;+\; 15 \cdot 10 \;+\; 12 \cdot 10
  \;=\; 180 + 150 + 120
  \;=\; 450.
\]

\begin{answerbox}
\centering
There are $\boxed{450}$ ways to choose two books of different subjects.
\end{answerbox}

% ══════════════════════════════════════════════════════════════════
\newpage
\section{Problem 5 --- Natural Numbers Less Than 700 Divisible by 5 and/or 3}
% ══════════════════════════════════════════════════════════════════

\begin{theorybox}[Counting Multiples in a Range]
The number of positive multiples of $d$ that are less than or equal to $N$ is
$\left\lfloor \frac{N}{d} \right\rfloor$.

More precisely, the multiples of $d$ in $\{1, 2, \ldots, N\}$ are
$d, 2d, 3d, \ldots, \lfloor N/d \rfloor \cdot d$,
so there are exactly $\lfloor N/d \rfloor$ of them.

\medskip
\textbf{``Less than $N$'' vs.\ ``at most $N$''.}\;
The natural numbers less than $N$ are $\{1, 2, \ldots, N-1\}$.
The number of multiples of $d$ among these is
$\left\lfloor \frac{N-1}{d} \right\rfloor$.
\end{theorybox}

\begin{problembox}
How many natural numbers less than 700 are divisible by 5?
How many of those are divisible by 3?
How many are divisible by both 3 and 5?
\end{problembox}

\subsection*{Solution}

We consider the set $S = \{1, 2, 3, \ldots, 699\}$ (natural numbers less than 700).

\paragraph{(a) Divisible by 5.}
The multiples of 5 in $S$ are $5, 10, 15, \ldots, 695$.
The number of such multiples is:
\[
  \left\lfloor \frac{699}{5} \right\rfloor = \lfloor 139.8 \rfloor = 139.
\]

\paragraph{(b) Divisible by 3 (among natural numbers less than 700).}
The multiples of 3 in $S$ are $3, 6, 9, \ldots, 699$.
The number of such multiples is:
\[
  \left\lfloor \frac{699}{3} \right\rfloor = 233.
\]

\paragraph{(c) Divisible by both 3 and 5.}
A number divisible by both 3 and 5 is divisible by $\mathrm{lcm}(3,5) = 15$.
The multiples of 15 in $S$ are $15, 30, 45, \ldots, 690$.
The number of such multiples is:
\[
  \left\lfloor \frac{699}{15} \right\rfloor = \lfloor 46.6 \rfloor = 46.
\]

\begin{answerbox}
\centering
Divisible by 5: $\boxed{139}$. \qquad
Divisible by 3: $\boxed{233}$. \qquad
Divisible by both 3 and 5: $\boxed{46}$.
\end{answerbox}

% ══════════════════════════════════════════════════════════════════
\newpage
\section{Problem 10 --- 7-Digit Numbers with Alternating Parity}
% ══════════════════════════════════════════════════════════════════

\begin{theorybox}[Product Rule for Digit Counting]
When counting numbers with digit-by-digit constraints, the product rule
is applied position by position.

\medskip
\textbf{Key consideration:}\;
The \emph{leading digit} (leftmost position) of an $n$-digit number
\textbf{cannot be 0}, since that would reduce the number of digits.
This restriction must always be accounted for separately.

\medskip
\textbf{Even digits:} $\{0, 2, 4, 6, 8\}$ --- five digits.\\
\textbf{Odd digits:} $\{1, 3, 5, 7, 9\}$ --- five digits.
\end{theorybox}

\begin{problembox}
Find the number of 7-digit numbers where the even-numbered positions
contain odd digits and the odd-numbered positions contain even digits.
\end{problembox}

\subsection*{Solution}

We label positions $1, 2, 3, 4, 5, 6, 7$ from left to right.

\begin{center}
\renewcommand{\arraystretch}{1.4}
\begin{tabular}{ccccccc}
\toprule
Pos.\ 1 & Pos.\ 2 & Pos.\ 3 & Pos.\ 4 & Pos.\ 5 & Pos.\ 6 & Pos.\ 7 \\
\midrule
\textit{odd} & \textit{even} & \textit{odd} & \textit{even} & \textit{odd} & \textit{even} & \textit{odd} \\
(position) & (position) & (position) & (position) & (position) & (position) & (position) \\
\midrule
even digit & odd digit & even digit & odd digit & even digit & odd digit & even digit \\
\bottomrule
\end{tabular}
\end{center}

\paragraph{Even positions (2, 4, 6):}
These must contain \textbf{odd} digits from $\{1, 3, 5, 7, 9\}$.
Each position has $5$ choices, giving $5^3$ combinations.

\paragraph{Odd positions (1, 3, 5, 7):}
These must contain \textbf{even} digits from $\{0, 2, 4, 6, 8\}$.
Each position has $5$ choices \emph{in general}, \textbf{except}
position 1 (the leading digit), which \textbf{cannot be 0}.

\begin{itemize}[nosep]
  \item Position 1: $\{2, 4, 6, 8\}$ --- \textbf{4 choices}.
  \item Positions 3, 5, 7: $\{0, 2, 4, 6, 8\}$ --- $5$ choices each, giving $5^3$.
\end{itemize}

\paragraph{Applying the product rule.}
Since the choices at each position are independent:
\[
  \underbrace{4}_{\text{pos.\ 1}}
  \;\times\;
  \underbrace{5^3}_{\text{pos.\ 2, 4, 6}}
  \;\times\;
  \underbrace{5^3}_{\text{pos.\ 3, 5, 7}}
  \;=\; 4 \times 125 \times 125
  \;=\; 4 \times 15{,}625
  \;=\; 62{,}500.
\]

\begin{answerbox}
\centering
There are $\boxed{62{,}500}$ such 7-digit numbers.
\end{answerbox}

% ══════════════════════════════════════════════════════════════════
\newpage
\part*{Section 2: Permutations and Arrangements}
\addcontentsline{toc}{part}{Section 2: Permutations and Arrangements}
% ══════════════════════════════════════════════════════════════════

% ══════════════════════════════════════════════════════════════════
\section{Problem 5 --- Boys and Girls in a Line}
% ══════════════════════════════════════════════════════════════════

\begin{theorybox}[Permutations with Positional Constraints]
When certain positions in a line are constrained (e.g., ``the first
and last positions must be occupied by specific types of people''),
we use the following strategy:

\begin{enumerate}[nosep]
  \item \textbf{Fill the constrained positions first},
        counting the number of ways to choose who goes there.
  \item \textbf{Fill the remaining positions} with the remaining people,
        who can stand in any order.
  \item \textbf{Multiply} the counts (product rule).
\end{enumerate}

\medskip
Recall that the number of ways to arrange $n$ distinct objects in a line
is $n! = n \cdot (n-1) \cdot (n-2) \cdots 2 \cdot 1$.
\end{theorybox}

\begin{problembox}
In how many ways can 6 boys and 7 girls stand in a line such that
the first and last positions are occupied by boys?
\end{problembox}

\subsection*{Solution}

There are $6 + 7 = 13$ people in total.

\paragraph{Step 1: Fill the constrained positions.}
\begin{itemize}[nosep]
  \item \textbf{First position:} must be a boy.
        There are 6 boys to choose from, so $6$ choices.
  \item \textbf{Last position:} must be a boy (different from the first).
        There are $6 - 1 = 5$ remaining boys, so $5$ choices.
\end{itemize}

\paragraph{Step 2: Fill the remaining positions.}
After placing two boys at the ends, there are $13 - 2 = 11$ people left
(4 boys and 7 girls). These 11 people can stand in the remaining
11 positions (positions 2 through 12) in any order:
\[
  11! \text{ ways.}
\]

\paragraph{Step 3: Apply the product rule.}
The choices at each step are independent, so:
\[
  6 \times 5 \times 11!
  \;=\; 30 \times 11!.
\]

\begin{answerbox}
\centering
The total number of arrangements is $\boxed{30 \times 11! = 30 \times 39{,}916{,}800 = 1{,}197{,}504{,}000}$.
\end{answerbox}

% ══════════════════════════════════════════════════════════════════
\newpage
\section{Problem 10 --- Ramsey's Theorem $R(3,3) = 6$}
% ══════════════════════════════════════════════════════════════════

\begin{theorybox}[Pigeonhole Principle and Ramsey Theory]
\textbf{Pigeonhole Principle.}\;
If $n$ objects are placed into $k$ boxes and $n > k$, then at least
one box contains more than one object.

More generally: if $n$ objects are placed into $k$ boxes, then at least
one box contains $\lceil n/k \rceil$ objects.

\medskip
\textbf{Ramsey Theory} (informal statement).
In any sufficiently large structure, order must appear.
The Ramsey number $R(s,t)$ is the smallest $n$ such that every
2-colouring (say, red and blue) of the edges of the complete graph
$K_n$ contains either a red $K_s$ or a blue $K_t$.

\medskip
The classical result $R(3,3) = 6$ states: in any group of 6 people,
there are always 3 mutual friends or 3 mutual strangers.
\end{theorybox}

\begin{problembox}
Prove that in any group of 6 people, there are either 3 mutual
friends or 3 mutual strangers.
\end{problembox}

\subsection*{Solution (Proof)}

We model the situation as a 2-colouring of the complete graph $K_6$.
Each person is a vertex, and each edge is coloured \textbf{red}
(friends) or \textbf{blue} (strangers).
We must show that there exists a monochromatic triangle.

\paragraph{Step 1: Apply the Pigeonhole Principle.}
Pick any vertex $A$. Vertex $A$ is connected to the other 5 vertices
by 5 edges, each coloured red or blue.
By the pigeonhole principle (5 edges into 2 colours),
at least $\lceil 5/2 \rceil = 3$ edges from $A$ have the same colour.

Without loss of generality, assume $A$ has at least 3 \textbf{red}
edges. Let the other endpoints of these red edges be $B$, $C$, $D$.
(The case with 3 blue edges is handled symmetrically by swapping
``friend'' and ``stranger.'')

\paragraph{Step 2: Examine the edges among $B$, $C$, $D$.}
Consider the three edges $BC$, $BD$, $CD$.
There are two cases:

\textbf{Case 1:} At least one of $BC$, $BD$, $CD$ is \textbf{red}.\\
Suppose (WLOG) $BC$ is red. Then $A$, $B$, $C$ form a red triangle:
edges $AB$, $AC$, and $BC$ are all red.
Thus we have 3 mutual \textbf{friends}.

\textbf{Case 2:} None of $BC$, $BD$, $CD$ is red, i.e., all three
are \textbf{blue}.\\
Then $B$, $C$, $D$ form a blue triangle: edges $BC$, $BD$, $CD$
are all blue.
Thus we have 3 mutual \textbf{strangers}.

\medskip
In both cases, a monochromatic triangle exists. \qed

\paragraph{Why 5 people are not enough.}
To show that $R(3,3) > 5$, we exhibit a 2-colouring of $K_5$
with no monochromatic triangle: arrange 5 vertices in a regular
pentagon and colour the edges of the pentagon red and the diagonals
blue (or vice versa). One can verify that no triangle is monochromatic.

\begin{answerbox}
\centering
In any group of 6 people, there are always 3 mutual friends
or 3 mutual strangers. This proves $R(3,3) = 6$.
\end{answerbox}

% ══════════════════════════════════════════════════════════════════
\newpage
\section{Problem 12 --- Divisibility Among $7, 77, 777, \ldots$}
% ══════════════════════════════════════════════════════════════════

\begin{theorybox}[Pigeonhole Principle for Divisibility]
A powerful application of the pigeonhole principle in number theory:

\medskip
If we have $n+1$ integers and consider their remainders modulo $n$,
then by the pigeonhole principle, at least two of them share the
same remainder. Their \emph{difference} is therefore divisible by $n$.

\medskip
\textbf{Repunit-like sequences.}\;
The sequence $7, 77, 777, 7777, \ldots$ can be written in closed form.
The $k$-th term is:
\[
  a_k = \underbrace{77\ldots7}_{k \text{ sevens}}
      = 7 \cdot \underbrace{11\ldots1}_{k}
      = 7 \cdot \frac{10^k - 1}{9}.
\]
\end{theorybox}

\begin{problembox}
Prove that among the first 2014 terms of the sequence
$7,\; 77,\; 777,\; 7777,\; \ldots$
there is one divisible by 2013.
\end{problembox}

\subsection*{Solution (Proof)}

Let $a_k = \underbrace{77\ldots7}_{k}$ for $k = 1, 2, \ldots, 2014$.

Consider the 2014 remainders $r_k = a_k \bmod 2013$ for $k = 1, 2, \ldots, 2014$.

\paragraph{Case 1: Some $r_k = 0$.}
If $r_k = 0$ for some $k$, then $a_k$ is divisible by 2013, and we are done.

\paragraph{Case 2: No $r_k = 0$.}
Then each $r_k \in \{1, 2, \ldots, 2012\}$.
We have 2014 remainders taking values in a set of size 2012.
By the \textbf{pigeonhole principle}, there exist indices $i < j$
(with $1 \le i < j \le 2014$) such that $r_i = r_j$, i.e.,
\[
  a_j \equiv a_i \pmod{2013}.
\]
Therefore $2013 \mid (a_j - a_i)$.

\paragraph{Analysing $a_j - a_i$.}
We compute:
\[
  a_j - a_i
  = \underbrace{77\ldots7}_{j} - \underbrace{77\ldots7}_{i}
  = \underbrace{77\ldots7}_{j-i} \cdot 10^i
  = a_{j-i} \cdot 10^i.
\]
So $2013 \mid a_{j-i} \cdot 10^i$.

Now, $2013 = 3 \times 11 \times 61$. Since $\gcd(10, 2013) = 1$
(as 2013 has no factors of 2 or 5), it follows that $\gcd(10^i, 2013) = 1$.

Therefore $2013 \mid a_{j-i}$.

Since $1 \le j - i \le 2013 < 2014$, the term $a_{j-i}$ is indeed
among the first 2014 terms of the sequence.

\medskip
In both cases, there exists a term in $\{a_1, a_2, \ldots, a_{2014}\}$
that is divisible by 2013. \qed

\begin{answerbox}
\centering
By the pigeonhole principle, among $a_1, a_2, \ldots, a_{2014}$,
there is at least one term divisible by $2013$.
\end{answerbox}

\end{document}
